\noindent326\hfill\textbf{BRIEFWECHSEL.\quad ZUR NETZAUSGLEICHUNG.}
\begin{center}\textsc{Gauss} an \textsc{Gerling}.\quad Göttingen, 14. November 1838.\end{center}

.\ .\ .\ .\ .\ .\ Gegen die Veröffentlichung der Formeln für die Anzahl der
Bedingungsgleichungen habe ich gar nichts einzuwenden. Es ist mir aber
\textls[120]{jetzt} unmöglich, die Ausnahmefälle präzise und erschöpfend zu
bezeichnen, und ich weiß auch nicht, ob dies überhaupt \textls[120]{möglich}
sein wird, ohne die Sache eingehender zu entwickeln. S[ine] m[editatione]
dürfte es genügen zu sagen, dass es in einem System von $p$ Punkten, zwischen
denen es $l$ Verbindungslinien gibt, wenn die Richtung jeder Verbindungslinie
an \textls[120]{beiden} Endpunkten gemessen, das heißt an eine andere Richtungslinie
des Systems geknüpft ist,

$l-p+1$ voneinander unabhängige Bedingungsgleichungen aus bloßen Winkelsummen
und außerdem noch

$l-2p+3$ \textls[120]{andere} Bedingungsgleichungen geben muss, die in den
\textls[120]{gewöhnlichsten} Fällen sämtlich unter die Kategorie der im
Supplem. Theor.\ p.\ 31 unter III[\textsuperscript{*})] angeführten fallen.
\par\medskip
\noindent\rule{.24\textwidth}{.4pt}\par\smallskip
{\footnotesize\noindent\hspace*{1.5em}[\textsuperscript{*}) Band IV, S.\ 81.]\par}
\vspace{2em}
\begin{center}\rule{.18\textwidth}{.4pt}\end{center}
